% =============================================================================
\section{Implementation}
\label{sec:implementation}
% =============================================================================

\subsection{System architecture}

The system is split into a Python generation server and a rendering client. The
server executes the pipeline and exposes the result as a \texttt{.frame} archive
over a WebSocket control channel and an HTTP asset channel, so a client streams
assets as they become available rather than waiting for the full run.

Two clients exist. A Unity client (C\#, URP) targets desktop VR and carries the
runtime shading work: a terrain splat shader blending the layered material of
\cref{sec:texture}, an animated water surface, and wind-driven foliage sway
sharing one global wind field so that grass, foliage and water gust together. A
native macOS client (Swift, Metal~4) renders the same scene and streams it to an
Apple Vision Pro through a remote immersive space.

Maintaining two renderers is a source of divergence worth recording: features
added to one do not automatically exist in the other. The water shader described
in \cref{sec:texture}, for instance, exists only in the Unity client, so on the
Metal client water falls back to the material parameters baked into the asset.
Keeping the baked material in step with the shader defaults is therefore not
cosmetic: it is what one of the two clients actually renders.

\subsection{Models}

\Cref{tab:models} lists the learned components. The pipeline is deliberately
model-agnostic where alternatives exist: panorama generation, object
reconstruction and inpainting each have several interchangeable backends selected
by configuration, which was essential during development because backend quality
varied substantially by scene type.

\begin{table}[t]
\centering
\small
\caption{Learned components and their role in the pipeline.}
\label{tab:models}
\begin{tabular}{@{}llp{6.1cm}@{}}
\toprule
\textbf{Role} & \textbf{Model} & \textbf{Notes} \\
\midrule
Captioning          & BLIP~\cite{li2022blipbootstrappinglanguageimagepretraining} & Scene and per-crop captions \\
Scene tagging       & RAM++~\cite{huang2023open,zhang2023recognize} & Open-set tags; used as a scene prior \\
Perspective depth   & Depth Anything 3~\cite{lin2025depth3recoveringvisual} & Input image \\
Panoramic depth     & DAP~\cite{lin2025dap} & Equirectangular-native \\
Panorama synthesis  & FLUX~\cite{esser2024scalingrectifiedflowtransformers} & Alternatives: DreamCube~\cite{dreamcube}, WorldGen~\cite{worldgen2025ziyangxie} \\
Super-resolution    & Swin2SR~\cite{conde2022swin2srswinv2transformercompressed} & Panorama upscaling \\
Segmentation        & SAM~2~\cite{ravi2024sam2segmentimages} & Objects; also video tracking \\
Matting             & BiRefNet~\cite{zheng2024birefnet} & High-resolution cutouts \\
Region typing       & SegFormer~\cite{DBLP:journals/corr/abs-2105-15203} & ADE20K classes folded into 9 region types \\
Classification      & CLIP & Fixed category set \\
Verification        & Grounding DINO~\cite{liu2024groundingdino} & Independent label corroboration \\
Similarity          & DINOv2 & Appearance clustering into buckets \\
Object 3D           & SAM~3D~\cite{sam3dteam2025sam3d3dfyimages} & Alternatives: SPAR3D~\cite{huang2025spar3dstablepointawarereconstruction}, TRELLIS~\cite{xiang2025trellis2} \\
Inpainting          & ObjectClear~\cite{zhao2026objectclear} & Alternatives: LaMa~\cite{suvorov2021resolution}, MAT~\cite{li2022mat} \\
Lighting            & LuxDiT~\cite{liang2025luxdit} & HDR environment estimation \\
De-lighting         & IntrinsicDiffusion~\cite{Luo2024IntrinsicDiffusion} & Albedo extraction for assets \\
Erosion             & Landlab~\cite{barnhart2020landlab} & Not learned; geomorphological simulation \\
\bottomrule
\end{tabular}
\end{table}

\subsection{Runtime and scene complexity}

A full run on the Mount Rainier capture takes \SI{93}{\minute} on a single GPU
across 38 active stages. The distribution is heavily skewed:
object segmentation (\SI{29.5}{\percent}) and per-group 3D reconstruction
(\SI{25.2}{\percent}) together account for over half, with video generation and
tracking (\SI{16.5}{\percent}) next. The entire terrain track (height field,
reconstruction, erosion, meshing and texturing) is under
\SI{10}{\percent}, which is worth noting given it is where most of this report's
technical contribution lies. Optimizing the pipeline means optimizing the learned
components, not the geometry.

The resulting scene contains a \num{79982}-triangle terrain mesh, a separate
\num{20228}-triangle collision mesh, 29 shared category meshes totalling
\SI{4.0}{\mega\byte}, and \num{9475} placed entities of which the overwhelming
majority are instanced ground cover. Because assets are shared per (class,
appearance) group and instanced, entity count and asset payload are effectively
decoupled: the property that makes a meadow of nine thousand plants affordable.
